% Generated draft: requires mathematical and source review.
\documentclass[11pt]{article}
\usepackage[margin=25mm]{geometry}
\usepackage{fontspec}
\setmainfont{Latin Modern Roman}
\usepackage{amsmath,amssymb,mathtools,unicode-math}
\setmathfont{Latin Modern Math}
\usepackage{enumitem,needspace,xurl,hyperref}
\setlength{\parindent}{0pt}
\setlength{\parskip}{5pt}
\setlength{\emergencystretch}{3em}
\newcommand{\N}{\mathbb{N}}
\newcommand{\Z}{\mathbb{Z}}
\newcommand{\Q}{\mathbb{Q}}
\newcommand{\R}{\mathbb{R}}
\newcommand{\C}{\mathbb{C}}
\newcommand{\F}{\mathbb{F}}
\newcommand{\A}{\mathbb{A}}
\newcommand{\PP}{\mathbb{P}}
\newcommand{\E}{\mathbb{E}}
\newcommand{\eps}{\varepsilon}
\newcommand{\dd}{\,\mathrm{d}}
\begin{document}
\section*{UnsolvedMath ID: 30001818}
\textbf{Source problem code:} catalogue entry
\subsection{Definitions and mathematical statement}
Let \(f\) be a normalized eigen-newform, let \(p\) be a prime, and let \(T\) denote the Galois representation attached to \(f\) in the supplied record. Put
\[
\mu_{p^\infty}=\bigcup_{n\geq 1}\mu_{p^n},\qquad
\Gamma=\operatorname{Gal}\bigl(\Q(\mu_{p^\infty})/\Q\bigr),\qquad
\Delta=\Gamma_{\mathrm{tors}},
\]
where \(\mu_{p^n}\) is the group of \(p^n\)-th roots of unity and \(\Delta\) is the torsion subgroup of \(\Gamma\). The record denotes the cyclotomic Iwasawa algebra by \(\Lambda\), the relevant cohomology modules by \(H^1(T)\) and \(H^2(T)\), and Kato's Euler-system element by \(z_f\). In the displayed quotient, \(\Lambda z_f\) is intended to be the submodule of \(H^1(T)\) generated by this element. The superscript \(\eta\) denotes the intended component associated with a character \(\eta\) of \(\Delta\), and \(\operatorname{Char}\) denotes a characteristic ideal over the corresponding Iwasawa coefficient ring.

The source-level question, attributed to Kato's main conjecture, is whether
\[
\operatorname{Char}\bigl(H^2(T)^\eta\bigr)
=
\operatorname{Char}\bigl((H^1(T)/\Lambda z_f)^\eta\bigr)
\qquad\text{for every character }\eta\text{ of }\Delta.
\]
Both characteristic ideals must be interpreted over the same appropriate coefficient ring, with the module hypotheses needed for their definition established.

\emph{Unresolved formulation qualifications.} The supplied record does not specify the weight, level, or character of \(f\), the allowed primes \(p\), the coefficient ring and lattice or twist conventions for \(T\), or the precise definition of \(\Lambda\). It also does not define the cohomology theory, inverse limits, ramification conditions, or local conditions underlying \(H^i(T)\). The normalization and coefficient realization of \(z_f\), the field of values of \(\eta\), and the construction of the \(\eta\)-component are unspecified. Finite generation and torsion of the modules whose characteristic ideals occur are not supplied. Consequently, this is the extracted conjectural equality with its universal quantifier over \(\eta\) preserved; the record alone does not determine a fully specified integral formulation or justify an assertion for every unrestricted choice of \(f\), \(p\), and \(T\).
\subsection{Short English statement}
The question asks whether, in the intended modular cyclotomic Iwasawa setting, the characteristic ideal of the second cohomology module equals that of the first cohomology module modulo Kato's Euler-system element, for every character of the finite torsion subgroup of the cyclotomic Galois group. Several essential conventions and hypotheses are missing from the supplied record, so the displayed equality requires further specification and human mathematical review.
\subsection{Sources}
\begin{itemize}
\item Supplied local UnsolvedMath registry record, permanent ID \(30001818\), title: Kato's Main Conjecture for Modular Galois Representations. Exact source problem code: \texttt{catalogue entry}. Collection: Oberwolfach Reports Open Problems. The record gives proposed year \(2011\) and names no proposer; the conjecture is attributed to Kato.
\item Supplied report citation: \emph{Algebraische Zahlentheorie} (2011), Oberwolfach Reports, pp.\ 1709--1768, DOI \texttt{10.4171/owr/2011/31}. Original report link: \url{https://doi.org/10.4171/owr/2011/31}. The supplied background describes Oberwolfach Reports as documenting mathematical workshops and their open-problem discussions; no specific page or contributor for this problem is supplied.
\item Supplied literature source, labelled ``Muller 2024'': \url{https://doi.org/10.1017/S0017089524000016}.
\item Supplied literature source, labelled ``Skinner--Urban'': \url{https://doi.org/10.1007/s00222-013-0448-1}.
\item Literature assessment dated \texttt{2026-08-21}: the supplied curation labels the problem partially solved and states that Kato's main conjecture is proved in important ordinary and CM cases, while the extracted broad formulation is not a theorem in full generality as of that check date. This is supplied dated literature triage derived from the OWR curation, not an independently verified assessment; its time-sensitive status and cited claims require independent verification.
\item The supplied registry labels the statement \texttt{corrected\_verified} and reports that its mathematical content was checked against the official Oberwolfach report, with notation, context, and question form normalized without changing the source-level claim. This is a supplied verification claim, not verification performed for this draft.
\end{itemize}
\end{document}
